% !TEX program = xelatex
\documentclass[aspectratio=169,11pt]{beamer}
% Fandol ships with common TeX distributions and avoids Windows-only font names.
\usepackage[UTF8,fontset=fandol]{ctex}
\setCJKfamilyfont{GYKai}[AutoFakeBold=2]{FandolKai-Regular}
\usepackage{amsmath,booktabs,array,tabularx}
\usepackage{GoatYangAcademic}
\AtBeginDocument{\CJKfamily{GYKai}}
\newcolumntype{Y}{>{\raggedright\arraybackslash}X}

% ===== Edit presentation metadata here =====
\title{GoatYang Academic}
\subtitle{学术汇报与组会 Beamer 模板}
\author[goat\_yang]{goat\_yang}
\institute{学院 / 研究机构 · 课题组}
\date{\today}
\shortreporttitle{GoatYang Academic}
% Optional user-provided logos; keep empty for the public example.
% \universitylogo{img/university-logo.pdf}
% \schoollogo{img/school-logo.png}
% \enableSectionPages
% \setbeamertemplate{navigation symbols}{}

\begin{document}
\begin{frame}[plain,noframenumbering]\titlepage\end{frame}
\begin{frame}{汇报目录}\tableofcontents[hideallsubsections]\end{frame}

\section{研究背景}
\begin{frame}[t]{研究背景与汇报主线}
\topic{核心问题}
\entry{用一句话说明研究对象、应用场景和当前方法的不足。}
\entry{明确本次汇报希望回答的问题，避免只罗列阅读过的材料。}
\contentdivider
\topic{汇报主线}
\entry{研究背景：为什么值得研究？}
\entry{方法分析：关键机制如何解决问题？}
\entry{验证计划：用什么实验支持结论？}
\end{frame}

\begin{frame}[t]{文献概览与选文依据}
\begin{tabularx}{\textwidth}{lYYY}
\toprule 编号 & 研究问题 & 核心方法 & 关注点 \\
\midrule
P1 & 基准问题与假设 & 基础方法 & 性能与成本 \\
P2 & 现有方法的不足 & 改进机制 & 消融证据 \\
P3 & 更广泛的场景 & 扩展方法 & 泛化与局限 \\
\bottomrule
\end{tabularx}
\vspace{.3cm}
\begin{block}{阅读记录}
填写真实题目、作者、年份与出处；将论文结论、个人理解和待验证假设分开。
\end{block}
\end{frame}

\section{方法分析}
\begin{frame}[t]{研究动机：问题与思路}
\begin{columns}[T,onlytextwidth]
\begin{column}{.48\textwidth}
\topic{现有问题}
\entry{基线方法在什么条件下表现不足？}
\entry{限制来自模型假设、搜索机制还是计算预算？}
\entry{已有实验证据能支持多强的判断？}
\end{column}
\begin{column}{.48\textwidth}
\topic{解决思路}
\entry{提出一个可以被独立验证的改进机制。}
\entry{解释机制与问题之间的对应关系。}
\entry{预先定义有效与无效的判断条件。}
\end{column}
\end{columns}
\end{frame}

\begin{frame}[t]{方法框架：从输入到验证}
\centering
\begin{tikzpicture}[node distance=.45cm,
  stage/.style={draw=GYBlue,fill=GYPaleBlue,rounded corners=3pt,
    minimum width=2.65cm,minimum height=1.15cm,align=center,font=\small},
  >=Stealth]
\node[stage] (a) {问题定义\\数据与约束};
\node[stage,right=of a] (b) {基线方法\\统一预算};
\node[stage,right=of b] (c) {改进机制\\独立消融};
\node[stage,right=of c] (d) {结果分析\\适用边界};
\draw[->,thick,GYBlue] (a)--(b);
\draw[->,thick,GYBlue] (b)--(c);
\draw[->,thick,GYBlue] (c)--(d);
\end{tikzpicture}
\par\vspace{.5cm}
\begin{block}{讲解重点}
先解释输入、输出与整体流程，再展开关键模块；每个模块对应一个明确的问题。
\end{block}
\end{frame}

\begin{frame}[t]{公式与符号说明}
\begin{columns}[T,onlytextwidth]
\begin{column}{.48\textwidth}
\topic{优化模型示例}
\begin{equation*}
\begin{aligned}
\min_{\boldsymbol{x}}\quad & f(\boldsymbol{x}) \\
\text{s.t.}\quad & g_j(\boldsymbol{x})\leq 0,\\
& \boldsymbol{\ell}\leq\boldsymbol{x}\leq\boldsymbol{u}.
\end{aligned}
\end{equation*}
\end{column}
\begin{column}{.48\textwidth}
\topic{符号含义}
\entry{$\boldsymbol{x}$：待优化的决策变量。}
\entry{$f$：目标函数；$g_j$：第 $j$ 个约束。}
\entry{$\boldsymbol{\ell},\boldsymbol{u}$：变量上下界。}
\entry{模型假设与单位需在实际汇报中补齐。}
\end{column}
\end{columns}
\end{frame}

\begin{frame}[t]{图文布局与图片尺寸}
\begin{columns}[T,onlytextwidth]
\begin{column}{.58\textwidth}
\centering
\begin{tikzpicture}
\node[draw=GYBlue,fill=GYPaleBlue,rounded corners,
 minimum width=.94\linewidth,minimum height=3.5cm,align=center]
 {在此放置方法图或实验图\\保持原始宽高比};
\end{tikzpicture}
\end{column}
\begin{column}{.38\textwidth}
\topic{图旁说明}
\entry{先说明读图方式，再给出主要发现。}
\entry{宽度和高度是容器上限，图片不会被拉伸。}
\entry{裁剪只去除无效空白，不切掉图例与坐标轴。}
\end{column}
\end{columns}
\end{frame}

\section{实验设计}
\begin{frame}[t]{最小验证与消融设计}
\begin{tabularx}{\textwidth}{lYY}
\toprule 组别 & 设置 & 验证目的 \\
\midrule
基线 & 原始方法与统一预算 & 建立可复现的参照 \\
改进 & 基线加单一机制 & 检查改进是否有效 \\
消融 & 移除关键组件 & 判断收益来源 \\
稳健性 & 多种场景与随机种子 & 判断结论适用范围 \\
\bottomrule
\end{tabularx}
\vspace{.3cm}
\begin{alertblock}{公平比较}
统一数据划分、评价预算、停止条件和统计口径，同时报告效果与计算成本。
\end{alertblock}
\end{frame}

\begin{frame}[t]{结果展示与结论边界}
\begin{columns}[T,onlytextwidth]
\begin{column}{.48\textwidth}
\begin{block}{观测结果}
填写实际实验指标、误差范围、样本量与运行条件。没有结果时保留为计划，不填虚构数字。
\end{block}
\end{column}
\begin{column}{.48\textwidth}
\begin{exampleblock}{解释与局限}
区分结果直接支持的结论与机制推测；列出失败场景、额外成本和待验证条件。
\end{exampleblock}
\end{column}
\end{columns}
\contentdivider
\topic{一句话结论}
本页回答一个具体问题，并用图表或数据支撑。
\end{frame}

\section{总结计划}
\begin{frame}[t]{下一步计划与讨论}
\topic{近期计划}
\entry{复现基线，核对数据与评价实现。}
\entry{完成单机制实验与必要消融。}
\entry{记录失败案例，调整研究假设。}
\contentdivider
\topic{需要讨论的问题}
\entry{研究问题是否清晰，验证范围是否合适？}
\entry{当前证据还缺少哪些关键对照？}
\end{frame}

\begin{frame}[plain,noframenumbering]
\centering\vfill
{\LARGE\bfseries\color{GYBlue}感谢聆听，欢迎讨论\par}
\vspace{.5cm}
{\color{GYGray}GoatYang Academic · goat\_yang\par}
\vfill
\end{frame}
\end{document}
